
We have modelled the labelling uncertainty and confusion probabilities in the
labelling process of the \textit{CIFAR-10H} dataset using a Bayesian latent
class mixture model. The resulting estimated confusion matrix was a very good
approximation of the true confusion matrix, thus validating the methodology. As
we have seen, the entropy can be used as a measure of the ambiguity of an image.
The prior entropy and the posterior entropy distributions of the images show
different characteristics, which can lead to different interpretations of what
makes an image ambiguous. We have also seen that the subcategories of the images
can differ significantly in their inherent image ambiguity. Furthermore, we have
analysed the time taken by the annotators to label the images and found that
a higher time taken to label an image is associated with a higher uncertainty in
the labelling process.\bigskip

This work can be used to improve the machine learning models in the future, by
incorporating the uncertainty in the labelling process into the training of the
models. This can be done by using the estimated confusion matrix as prior
information in the training of the models. This can lead to more robust models
that are less sensitive to noisy labels. The fact that we built the model based
on a machine learning benchmark dataset, the \textit{CIFAR-10} dataset, makes
the results of this work applicable to a wide range of machine learning models.
The methodology can be extended to other datasets and other labelling processes,
where the true labels are unknown.\bigskip

There are some limitations though. The model is based on the assumption that
there is a ground truth, i.e. a true labels, which is unknown. This assumption
may not hold in some cases, where there is just no ground truth. In such cases,
the model may not be applicable. Furthermore, the model is based on the
assumption that the images are independent of each other. This may not hold in
some cases, where the images are correlated. In such cases, the model may not be
applicable. Another problematic aspect that might occur are numerical
instabilities or computational limitations in the estimation of the model
parameters, which can lead to unreliable results. This can result from too many
latent classes or very small confusion probabilities or too many votes per
image. In such cases, we have to be careful about the correctness of the
results. Lastly, the model is based on the assumption that the confusion
probabilities are constant over time. This may not hold in some cases, where the
confusion probabilities change over time, due to the annotators getting better
or worse at labelling the images or due to the images themselves changing over
time.\bigskip

For future research, it would be interesting to investigate the effect of
labelling uncertainty on the performance of machine learning models. This can be
done by comparing the performance of models trained with and without the prior
information of the estimated confusion matrix. We could extend the work of
\cite{kollerGoingOneHotEncoding2024} to include the reaction time of the
annotators as a feature in the model and apply their methodology to the
\textit{CIFAR-10H} dataset. We could also apply our methodology to other
datasets, such as the \textit{ImageNet} dataset, and compare the results to those
obtained with the \textit{CIFAR-10H} dataset, or apply it to datasets
from other domains, such as medical imaging or quality control in manufacturing.